\documentclass[11pt]{article}
\usepackage[margin=1in]{geometry}
\usepackage{amsmath}
\usepackage{amssymb}
\usepackage{amsfonts}

\title{A Brief Review of Core Concepts in Statistical Learning}
\author{J.\ Doe \and A.\ Smith}
\date{}

\begin{document}
\maketitle

\begin{abstract}
We review several foundational quantities that recur throughout statistical
learning: information divergences, the softmax output layer, empirical risk
minimization via stochastic gradient updates, mutual information, and the
Gaussian likelihood. The exposition is self-contained and standalone display
equations are marked in the usual way.
\end{abstract}

\section{Information Divergences}

The Kullback--Leibler divergence between two continuous probability
distributions $P$ and $Q$ with densities $p$ and $q$ is given by

\begin{equation}
D_{\mathrm{KL}}(P \| Q) = \int_{-\infty}^{\infty} p(x) \ln \frac{p(x)}{q(x)} \, dx.
\end{equation}

It is non-negative and vanishes if and only if $P = Q$ almost everywhere. When
comparing discrete categorical outputs, it is common to first squash logits
$\mathbf{z} \in \mathbb{R}^K$ through the softmax function

\begin{equation}
\sigma(\mathbf{z})_k = \frac{e^{z_k}}{\sum_{j=1}^{K} e^{z_j}},
\end{equation}

which yields a point on the probability simplex.

\section{Empirical Risk and SGD}

For binary classification the negative log-likelihood reduces to the familiar
cross-entropy loss,

\begin{equation}
\mathcal{L}(\theta) = -\frac{1}{N} \sum_{i=1}^{N} \left\{ y_i \log \hat{y}_i + (1 - y_i) \log(1 - \hat{y}_i) \right].
\end{equation}

Training then proceeds by stochastic gradient descent on the expected loss,

\begin{equation}
\theta_{t+1} = \theta_t - \eta \nabla_\theta \mathbb{E}_{(x,y) \sim \mathcal{D}} \left[ \ell(f_\theta(x), y) \right],
\end{equation}

where $\eta > 0$ is a step size and $\mathcal{D}$ is the data distribution.

\section{Mutual Information}

For a pair of discrete random variables $X$ and $Y$ with joint distribution
$p(x,y)$, the mutual information quantifies the reduction in uncertainty of
one when the other is observed:

\begin{equation}
I(X; Y) = \sum_{x \in \mathcal{X}} \sum_{y \in \mathcal{Y}} p(x, y) \log \frac{p(x, y)}{p(x)\, p(y)}.
\end{equation}

\section{The Gaussian Likelihood}

Finally, many continuous models assume Gaussian residuals. The density of a
univariate normal with mean $\mu$ and variance $\sigma^2$ is

\begin{equation}
f(x \mid \mu, \sigma^2) = \frac{1}{\sqrt{2\pi\sigma^2}} \exp\!\left( -\frac{(x - \mu)^2}{2\sigma^2} \right).
\end{equation}

Combined with i.i.d.\ assumptions this supports both maximum-likelihood
estimation and Bayesian posterior inference.

\end{document}
